\documentclass[11pt]{article}

\usepackage{nmrdoc}
\usepackage{siunitx}
\usepackage{bm}
\usepackage{tikz}

\sisetup{per-mode=symbol}
\DeclareSIUnit{\angstrom}{\text{\AA}}
\DeclareSIUnit{\ppm}{ppm}

\newcommand{\code}[1]{\texttt{#1}}
\setlength{\parindent}{0pt}
\setlength{\parskip}{0.42em}
\emergencystretch=2em

\begin{document}
\NmrTitle{Numerical Addendum: LarsenHBond Calculator}

The LarsenHBond calculator is an empirical tensor-grid lookup followed by a
frame change and a spherical-tensor decomposition. The lookup key is a donor
class, an acceptor class, and the three geometry coordinates
\((r,\theta,\rho)\). LarsenHBond's output is already in \(\si{\ppm}\): the
grid entries are DFT shielding-difference tensors from the Larsen
parameterization, not a geometric kernel awaiting calibration. The listings
below are distilled from the source, with loops and bookkeeping trimmed, but
the constants, guards, inequalities, and multiplication order are the
implementation's.

\section{The lookup point}

The donor sweep first makes a finite candidate set. Amide hydrogens and alpha
hydrogens are donor candidates; acceptor oxygen candidates come from the atom
spatial index within \(\SI{4.2}{\angstrom}\) of the donor hydrogen. The
acceptor classifier supplies the acceptor class and the two acceptor-side
atoms used below. For a classified candidate, the code computes
\((r,\theta,\rho)\), rejects \(\theta<\ang{90}\) and
\(\theta>\ang{180}\), then asks the grid for an interpolated record. Equality
at \(\ang{90}\) and \(\ang{180}\) reaches the grid lookup.

\begin{codebox}
constexpr double kSpatialCutoff_A = 4.2;
constexpr double kThetaMinDeg     = 90.0;
constexpr double kThetaMaxDeg     = 180.0;

auto candidate_atoms =
    spatial.AtomsWithinRadius(donor_pos, kSpatialCutoff_A);

LarsenHBondGeometry geom = ComputeLarsenHBondGeometry(
    donor_H_pos, accept_O_pos, accept_C_pos, accept_thd_pos);

if (geom.theta_deg < kThetaMinDeg || geom.theta_deg > kThetaMaxDeg)
    return PairResult::ThetaOutOfRange;

LarsenHBondRecord rec = grid.QueryNearest(donor_class, acc.class_, geom);
if (!rec.IsHit()) return PairResult::GridMiss;
\end{codebox}

The spatial cutoff, the two theta limits, and the water offset
\(\Delta\sigma_w=\SI{2.07}{\ppm}\) (the isotropic term added later to amide
hydrogens for which the sweep finds no \(\theta\)-accepted geometric pair) are
compiled constants in this calculator.
The H-bond keys in \code{data/calculator\_params.toml} belong to the separate
kernel-form hydrogen-bond calculators and are not read by this grid lookup.

In the geometry helper, \(\bm h\) is the donor hydrogen, \(\bm o\) the
acceptor oxygen, \(\bm c\) the heavy atom bonded to the acceptor oxygen, and
\(\bm t\) the acceptor-side third atom. The distance is
\(r=\|\bm o-\bm h\|\). The angle \(\theta\) is the H--O--C angle at the
acceptor oxygen. The cosine is clamped to \([-1,1]\) before \(\arccos\), so a
rounding excursion at the dot-product boundary becomes the endpoint angle
rather than a NaN.

The dihedral \(\rho\) compares the H--O--C plane to the O--C--third plane. The
two cross products make the plane normals, and the \(\operatorname{atan2}\)
pair keeps the sign and quadrant. The value is in degrees and is not negated
before lookup; the dense grids are keyed on this same sign convention.

\begin{codebox}
Vec3 H_to_O = acceptor_O_pos - donor_H_pos;
geom.r_angstrom = H_to_O.norm();

Vec3 O_to_H = donor_H_pos    - acceptor_O_pos;
Vec3 O_to_C = acceptor_C_pos - acceptor_O_pos;
double cos_theta = O_to_H.dot(O_to_C) /
                   (O_to_H.norm() * O_to_C.norm());
cos_theta = std::clamp(cos_theta, -1.0, 1.0);
geom.theta_deg = std::acos(cos_theta) * (180.0 / M_PI);

Vec3 C_to_third = acceptor_third_pos - acceptor_C_pos;
Vec3 n1 = H_to_O.cross(O_to_C);             // plane H-O-C
Vec3 n2 = O_to_C.cross(C_to_third);         // plane O-C-third
Vec3 m1 = n1.cross(O_to_C.normalized());
double x = n1.dot(n2);
double y = m1.dot(n2);
geom.rho_deg = std::atan2(y, x) * (180.0 / M_PI);
\end{codebox}

The donor and acceptor classes select one of six archives. Side-chain
carbonyl acceptors fold into the backbone-carbonyl archive because the Larsen
grid set has no separate side-chain-amide acceptor scan. In accumulation, the
acceptor-side secondary readouts are routed only for backbone-carbonyl
acceptors with a valid downstream residue. The archive choice selects the
tensor arrays; the \((r,\theta,\rho)\) coordinates select the cell inside
that archive.

\begin{codebox}
bool nma_acceptor =
    acceptor == HBondAcceptorClass::BackboneCarbonyl ||
    acceptor == HBondAcceptorClass::SidechainCarbonyl;

if (donor == HBondDonorClass::AmideHydrogen) {
    if (nma_acceptor)                                      return kArchiveNMANMA;
    if (acceptor == HBondAcceptorClass::HydroxylOxygen)    return kArchiveNMACOH;
    if (acceptor == HBondAcceptorClass::CarboxylateOxygen) return kArchiveNMACOO;
} else {
    if (nma_acceptor)                                      return kArchiveALANMA;
    if (acceptor == HBondAcceptorClass::HydroxylOxygen)    return kArchiveALACOH;
    if (acceptor == HBondAcceptorClass::CarboxylateOxygen) return kArchiveALACOO;
}
\end{codebox}

\section{Trilinear tensor interpolation}

Each readout dataset is a flat row-major array with logical shape
\((N_r,N_\theta,N_\rho,3,3)\). A successful query first locates a lower index
and a fractional coordinate on each axis. The \(r\) and \(\theta\) axes are
bounded axes. The tolerance is \(10^{-9}\) of one axis step: values below the
lower bound by more than that, or above the upper bound by more than that,
return a non-hit. Values inside that tolerance are clamped onto the stored
bound, so the lower bound has fraction \(0\) and the upper bound has fraction
\(1\) in the last cell.

The \(\rho\) axis is different. It is periodic, with a required
\(\ang{360}\) period. The record first wraps \(\rho\) into
\([-\ang{180},\ang{180})\); the periodic axis lookup then locates the cell in
the half-open period. In the current dense grids the stored \(\rho\) axis
runs from \(-\ang{180}\) to \(\ang{175}\) in \(\ang{5}\) steps, so the upper
corner of the last cell wraps to index \(0\). A value near
\(\ang{180}\) blends the \(\ang{175}\) plane with the \(-\ang{180}\) plane;
\(\ang{180}\) itself wraps to \(-\ang{180}\).

\begin{codebox}
rec.rho_deg = WrapRho(geom.rho_deg);

AxisLookup lr =
    LookupAxis(g.r_axis,     geom.r_angstrom, false);
AxisLookup lth =
    LookupAxis(g.theta_axis, geom.theta_deg,  false);
AxisLookup lrho =
    LookupAxis(g.rho_axis,   rec.rho_deg,     true);

if (lr.idx < 0 || lth.idx < 0 || lrho.idx < 0)
    return rec;                              // non-hit
\end{codebox}

\begin{codebox}
// Non-periodic path, distilled from LookupAxis.
double step = axis[1] - axis[0];
double tol  = std::abs(step) * kAxisBoundTolerance;  // 1e-9
if (value < axis.front() - tol || value > axis.back() + tol)
    return miss;

double v = std::clamp(value, axis.front(), axis.back());
double f = (v - axis.front()) / step;
int i = static_cast<int>(std::floor(f));
i = std::min(i, n - 2);
i = std::max(i, 0);

out.idx      = i;
out.idx_next = i + 1;
out.frac     = std::clamp(f - i, 0.0, 1.0);
\end{codebox}

\begin{codebox}
// Periodic path, distilled from LookupAxis.
double step   = axis[1] - axis[0];
double period = axis.back() - axis[0] + step;
if (std::abs(period - 360.0) > 1e-6) return miss;

double w = axis[0] + std::fmod(value - axis[0], period);
if (w < axis[0]) w += period;

double f = (w - axis[0]) / step;
int i = static_cast<int>(std::floor(f));
int i_next = i + 1;
if (i_next >= n) i_next = 0;       // last rho cell wraps

out.idx      = i;
out.idx_next = i_next;
out.frac     = f - i;
\end{codebox}

The interpolation itself is done independently for each matrix entry. For a
fixed row and column, the code reads the eight corner values
\(c_{\alpha\beta\gamma}\), blends first along \(\rho\), then along
\(\theta\), then along \(r\). This is the same product-weight blend as
\[
  \sigma_{ab}(r,\theta,\rho) =
  \sum_{\alpha,\beta,\gamma\in\{0,1\}}
  (1-f_r)^{1-\alpha} f_r^\alpha
  (1-f_\theta)^{1-\beta} f_\theta^\beta
  (1-f_\rho)^{1-\gamma} f_\rho^\gamma
  G_{i+\alpha,j+\beta,k+\gamma,a,b}.
\]
The eight weights sum to one for each matrix entry. The dense HDF5 grids
contain the cubic-spline pre-computation; the runtime path is trilinear.

\begin{codebox}
// Distilled from TrilinearMat3, one row/col entry at a time.
double c000 = cell(ir,      ith,      irho);
double c001 = cell(ir,      ith,      irho_next);
double c010 = cell(ir,      ith_next, irho);
double c011 = cell(ir,      ith_next, irho_next);
double c100 = cell(ir_next, ith,      irho);
double c101 = cell(ir_next, ith,      irho_next);
double c110 = cell(ir_next, ith_next, irho);
double c111 = cell(ir_next, ith_next, irho_next);

double c00 = c000 * (1.0 - frho) + c001 * frho;
double c01 = c010 * (1.0 - frho) + c011 * frho;
double c10 = c100 * (1.0 - frho) + c101 * frho;
double c11 = c110 * (1.0 - frho) + c111 * frho;

double c0 = c00 * (1.0 - fth) + c01 * fth;
double c1 = c10 * (1.0 - fth) + c11 * fth;

out(row, col) = c0 * (1.0 - fr) + c1 * fr;
\end{codebox}

When a validity mask is present, the same eight corners are inspected. If any
corner has mask value zero, \code{rec.any\_corner\_imputed} is set. The flag
does not change the interpolation weights; it records that the returned tensor
used at least one dense cell filled upstream for continuity.

\section{The donor frame}

The grid tensors are stored in a donor-centered canonical frame. The runtime
rebuilds that frame from the protein-side donor atoms before rotating any
looked-up tensor. For an amide donor, the anchor is the amide N and the third
atom is the predecessor residue's carbonyl C. For an alpha donor, the anchor
is \(\mathrm{C\alpha}\) and the third atom is the same residue's N.

The rotation matrix \(R\) stores the canonical basis vectors as rows in the
protein lab coordinates. The \(z\)-axis is fixed first:
\(\hat{\bm z}=(\bm h-\bm a)/\|\bm h-\bm a\|\). Then
\(\bm q=\bm h-\bm u\), with \(\bm u\) the donor-side third atom, loses its
\(\hat{\bm z}\) component:
\[
  \bm x_{\mathrm{raw}} =
  \bm q - (\bm q\cdot\hat{\bm z})\hat{\bm z}.
\]
Normalizing that projected vector gives \(\hat{\bm x}\), and
\(\hat{\bm y}=\hat{\bm z}\times\hat{\bm x}\). The three guards test norms
with a strict \(<10^{-9}\) comparison. If a guard fires, the helper returns
the identity matrix and logs a warning; the caller then uses that identity in
the same rotation path as any other \(R\).

\begin{codebox}
constexpr double kTinyVec = 1e-9;

Vec3 z_raw = donor_H_pos - donor_anchor_pos;
if (z_raw.norm() < kTinyVec) return Mat3::Identity();
Vec3 z = z_raw.normalized();

Vec3 third_to_H = donor_H_pos - donor_third_pos;
if (third_to_H.norm() < kTinyVec) return Mat3::Identity();

Vec3 x_raw = third_to_H - (third_to_H.dot(z)) * z;
if (x_raw.norm() < kTinyVec) return Mat3::Identity();

Vec3 x = x_raw.normalized();
Vec3 y = z.cross(x);

Mat3 R;
R.row(0) = x;
R.row(1) = y;
R.row(2) = z;
\end{codebox}

\begin{figure}[ht]
\centering
\begin{tikzpicture}[scale=1.05,line width=0.45pt,>=stealth]
  \coordinate (H) at (0,0);
  \coordinate (A) at (0,-2.0);
  \coordinate (U) at (-1.8,-0.75);
  \coordinate (P) at (0,-0.75);
  \coordinate (X) at (1.55,0);
  \coordinate (Z) at (0,1.45);

  \draw[black!45] (A) -- (H) -- (U);
  \draw[densely dashed] (U) -- (H)
      node[midway,above left] {\scriptsize \(\bm q=\bm h-\bm u\)};
  \draw[dotted] (U) -- (P);
  \draw[dotted] (P) -- (H);

  \draw[->] (H) -- (Z) node[above] {\scriptsize \(\hat{\bm z}\)};
  \draw[->] (H) -- (X) node[right] {\scriptsize \(\hat{\bm x}\)};
  \draw (1.0,0.0) circle (0.13);
  \fill (1.0,0.0) circle (0.035);
  \node at (1.32,0.20) {\scriptsize \(\hat{\bm y}=\hat{\bm z}\times\hat{\bm x}\)};

  \fill (H) circle (1.5pt) node[below right] {\scriptsize donor H \(\bm h\)};
  \fill (A) circle (1.5pt) node[below] {\scriptsize anchor \(\bm a\)};
  \fill (U) circle (1.5pt) node[left] {\scriptsize third \(\bm u\)};

\end{tikzpicture}
\caption{The donor frame keeps the anchor-to-H direction as
\(\hat{\bm z}\). The vector \(\bm h-\bm u\) supplies the second direction; one
projection removes its \(\hat{\bm z}\) component before normalization gives
\(\hat{\bm x}\). The \(\hat{\bm y}\) axis is the cross product
\(\hat{\bm z}\times\hat{\bm x}\).}
\end{figure}

\section{Rotating the tensor}

The stored-frame convention fixes the multiplication direction. During grid
generation, each log tensor was stored as
\[
  \bm\sigma_{\mathrm{canon}} =
  R_{\mathrm{log}}\,
  \bm\sigma_{\mathrm{log}}\,
  R_{\mathrm{log}}^{\mathsf T}.
\]
The protein-side \(R_{\mathrm p}\) has the same row convention, so recovering
the protein lab-frame tensor puts the transpose on the left:
\[
  \bm\sigma_{\mathrm{lab}} =
  R_{\mathrm p}^{\mathsf T}\,
  \bm\sigma_{\mathrm{canon}}\,
  R_{\mathrm p}.
\]
Both indices of the \(3\times3\) shielding tensor move through the same frame
change. The code applies this helper to donor-side and acceptor-side readouts;
the frame is still the donor frame for both, because the archive stores every
readout in that canonical donor basis.

\begin{codebox}
Mat3 RotateTensorToProteinLabFrame(
    const Mat3& sigma_canonical,
    const Mat3& R_protein) {
    return R_protein.transpose() * sigma_canonical * R_protein;
}

Mat3 R_protein = ComputeLarsenDonorFrame(
    donor_H_pos, donor_anchor, donor_third);

Mat3 sigma_lab = RotateTensorToProteinLabFrame(
    readout.canonical_tensor, R_protein);

AccumulateContribution(conf.MutableAtomAt(target_ai),
                       term, sigma_lab);
\end{codebox}

The four pair-term accumulators are \(3\times3\) sums in
\(\si{\ppm}\). At the end of the conformation, the total pair tensor for an
atom is the sum of the primary and secondary amide-H terms and the primary and
secondary alpha-H terms. The water term is a separate scalar field on amide H
atoms; it is not added to the tensor total.

\begin{codebox}
Mat3 shielding = a.larsen_hbond_1pHB_tensor
               + a.larsen_hbond_2pHB_tensor
               + a.larsen_hbond_1pHaB_tensor
               + a.larsen_hbond_2pHaB_tensor;

a.larsen_hbond_shielding_tensor    = shielding;
a.larsen_hbond_shielding_spherical =
    SphericalTensor::Decompose(shielding);
\end{codebox}

\section{The shared tensor split}

Every nine-column Larsen shielding array written by \code{WriteFeatures()} is
produced by the shared decomposition below. That function is closed-form
arithmetic on the matrix as accumulated; it does not pre-symmetrize or
diagonalize the matrix. The trace average gives \(T_0\), the antisymmetric
part gives \(T_1\), and the symmetric traceless part gives the five \(T_2\)
numbers in the shared packed order.

\begin{codebox}
st.T0 = sigma.trace() / 3.0;

st.T1[0] = 0.5 * (sigma(1,2) - sigma(2,1));
st.T1[1] = 0.5 * (sigma(2,0) - sigma(0,2));
st.T1[2] = 0.5 * (sigma(0,1) - sigma(1,0));

double Sxx = sigma(0,0) - st.T0;
double Syy = sigma(1,1) - st.T0;
double Szz = sigma(2,2) - st.T0;
double Sxy = 0.5 * (sigma(0,1) + sigma(1,0));
double Sxz = 0.5 * (sigma(0,2) + sigma(2,0));
double Syz = 0.5 * (sigma(1,2) + sigma(2,1));

st.T2[0] = std::sqrt(2.0) * Sxy;
st.T2[1] = std::sqrt(2.0) * Syz;
st.T2[2] = std::sqrt(3.0 / 2.0) * Szz;
st.T2[3] = std::sqrt(2.0) * Sxz;
st.T2[4] = (Sxx - Syy) / std::sqrt(2.0);
\end{codebox}

The \(T_2\) normalization is isometric. An off-diagonal entry such as
\(S_{xy}\) appears twice in the symmetric matrix, so it contributes
\(2S_{xy}^2\) to the Frobenius norm. The stored value
\(\sqrt2\,S_{xy}\) squares to that same amount. The two diagonal entries
\(\sqrt{3/2}\,S_{zz}\) and \((S_{xx}-S_{yy})/\sqrt2\) carry the two diagonal
degrees of freedom left after \(S_{xx}+S_{yy}+S_{zz}=0\). Thus a dot product
of two stored \(T_2\) five-vectors is the Frobenius inner product of the
symmetric traceless matrices they encode.

For LarsenHBond, the grid can return full \(3\times3\)
shielding-difference tensors, so all three channels can be populated.
\code{WriteFeatures()} emits six \code{SphericalTensor}-packed float64 arrays
of shape \(N\times9\): the total pair shielding, the four pair-term
shieldings, and the \(\mathrm{C\beta}\) diagnostic shielding. It also emits
the scalar water term as a length-\(N\) float64 vector in \(\si{\ppm}\) and
the per-atom pair count as a length-\(N\) int32 vector.

\section{Glossary}

\textbf{Trilinear interpolation.} A point inside a grid cell borrows from its
eight surrounding corners, the nearer corners weighing more (the eight corner
values blended by the point's three fractional distances). For one Larsen tensor
entry, the code blends the two \(\rho\) corners, blends
those results along \(\theta\), then blends the two remaining values along
\(r\). Formally, it is the tensor-product linear interpolant on a rectangular
cell, with weights
\((1-f_r)^{1-\alpha}f_r^\alpha
 (1-f_\theta)^{1-\beta}f_\theta^\beta
 (1-f_\rho)^{1-\gamma}f_\rho^\gamma\).

\textbf{Periodic lookup axis.} A ruler bent into a ring, so the cell past its
last mark uses the first mark as its upper corner. In the Larsen
\(\rho\) grid, a value near \(\ang{180}\) uses the cell from \(\ang{175}\) to
\(-\ang{180}\), while \(r\) and \(\theta\) outside their tolerated bounds
return a non-hit. Formally, the coordinate is reduced modulo the period before
the lower index and fractional coordinate are computed.

\textbf{Gram--Schmidt donor frame.} Three atoms raise a set of mutually
perpendicular axes --- one bond fixes the first, a second bond is straightened
perpendicular to it to set the next, and their cross product gives the third (an
orthonormal frame: the \(z\)-axis, then the projected \(x\)-axis, then
\(\hat{\bm y}=\hat{\bm z}\times\hat{\bm x}\)). Here \(\bm h-\bm a\) fixes
\(\hat{\bm z}\), and the component of
\(\bm h-\bm u\) perpendicular to \(\hat{\bm z}\) fixes \(\hat{\bm x}\).
Formally, this is one Gram--Schmidt projection followed by
\(\hat{\bm y}=\hat{\bm z}\times\hat{\bm x}\), with identity returned on the
three guarded degeneracies.

\textbf{Signed dihedral.} A signed dihedral measures how one plane turns into
the next around their shared bond. In the H--O--C--third coordinate, the
normal to H--O--C and the normal to O--C--third are compared with
\(\operatorname{atan2}\), so crossing the reference plane changes the sign of
\(\rho\). Formally, with
\(\bm n_1=(\bm o-\bm h)\times(\bm c-\bm o)\) and
\(\bm n_2=(\bm c-\bm o)\times(\bm t-\bm c)\),
\(\rho=\operatorname{atan2}((\bm n_1\times\widehat{\bm c-\bm o})\cdot\bm n_2,
\bm n_1\cdot\bm n_2)\).

\textbf{Spherical tensor split.} A \(3\times3\) tensor holds three kinds of
content that rotate without mixing --- an orientation-free size (the scalar
trace, \(T_0\)), an axial twist from its antisymmetric part (the three \(T_1\)
components), and a five-number traceless shape for the rest of the
direction-dependence (the five \(T_2\) components). A Larsen grid tensor can
carry an isotropic shielding difference in \(T_0\), an antisymmetric residue in
\(T_1\), and anisotropic shielding in \(T_2\). Formally, it is the closed-form
split of a \(3\times3\) tensor's nine components into dimensions \(1+3+5\), with
the \(T_2\) basis normalized to preserve the Frobenius norm.

\end{document}
